\subsection{道路的提升}

定理 3. --- 设 I 是 R 的一个区间，X 是一个非空分离拓扑空间。设 p: X → I 是连续、开且逆紧的映射，其纤维全不连通（TG, I, p.~83）。映射 p 是满射。对 X 中的每个点 x，存在 p 的一个连续截面 s，使得 s(p(x)) = x。

集合 \(p(\mathrm{X})\) 是 I 的一个开、闭（TG, I, p.~72, 命题 1）且非空的子集。由于 I 是连通的，有 \(p(\mathrm{X}) = \mathrm{I}\)；因此映射 \(p\) 是满射。

对每个 \((a,b)\in\mathrm{I}\times\mathrm{I}\) 满足 \(a\leqslant b\) 的，记 \(F_{a,b}\) 为由 \((y,z)\in\overline{p}^{-1}(a)\times\overline{p}^{-1}(b)\) 组成的集合，其中 y 和 z 属于 \(\overline{p}^{-1}([a,b])\) 的同一个连通分支。

引理 10. --- 设 \(a\) 和 \(b\) 是 I 中满足 \(a \leqslant b\) 的点。

\begin{enumerate}
\def\labelenumi{\alph{enumi})}
\item
  \(F_{a,b}\) 在 \(\overline{p}^{-1}(a)\times\overline{p}^{-1}(b)\) 中是闭的。
\item
  有 \(\operatorname{pr}_1(\mathrm{F}_{a,b}) = \overline{p}^{-1}(a)\) 且 \(\operatorname{pr}_2(\mathrm{F}_{a,b}) = \overline{p}^{-1}(b)\)。
\item
  设 \(c \in I\) 满足 \(b \leqslant c\)。若 \((y,z)\) 属于 \(F_{a,b}\) 且 \((z,t)\) 属于 \(F_{b,c}\)，则 \((y,t)\) 属于 \(F_{a,c}\)。
\end{enumerate}

集合 \(\overline{p}^{-1}([a,b])\) 是紧的（TG, I, p.~77, 命题 7）。因此，\((y,z)\in\overline{p}^{-1}(a)\times\overline{p}^{-1}(b)\) 属于 \(F_{a,b}\) 的充要条件是：\(\overline{p}^{-1}([a,b])\) 中每个包含 \(y\) 的开闭子集也包含 \(z\)（TG, II, p.~32, 命题 6）。

置 \(\mathrm{Y} = \overline{p}^{-1}([a,b])\)。对于 Y 的每个既开又闭的子集 U，集合 \(((\mathrm{U} \times \mathrm{U}) \cup ((\mathrm{Y} - \mathrm{U}) \times (\mathrm{Y} - \mathrm{U}))) \cap (\overline{p}^{-1}(a) \times \overline{p}^{-1}(b))\) 在 \(\overline{p}^{-1}(a) \times \overline{p}^{-1}(b)\) 中是闭的。这些集合的交等于 \(\mathrm{F}_{a,b}\)，故得出 \(a\)。

令 \(y \in \overline{p}^{-1}(a)\)。记 \(\mathcal{U}\) 为 Y 中 y 的开闭邻域组成的集合。由 p 取子空间诱导的从 Y 到 {[}a, b{]} 的映射是开且逆紧的（I, p.~17, 命题 8），因而也是闭的。于是，对每个 \(U\) 属于 \(\mathcal{U}\) 的，\(p(U)\) 是 {[}a, b{]} 的非空开闭子集；由于区间 {[}a, b{]} 是连通的，有 \(p(U) = [a, b]\)，特别地 \(U \cap \overline{p}^{-1}(b) \neq \varnothing\)。因此，\((U \cap \overline{p}^{-1}(b))_{U \in \mathcal{U}}\) 是紧空间 \(\overline{p}^{-1}(b)\) 的一个由非空闭子集组成的递减过滤族。该族的交非空（TG, I, p.~59）；令 z 为其中一个元素。有 \((y, z) \in \mathrm{F}_{a,b}\)。我们证明了关系 \(\mathrm{pr}_{1}(\mathrm{F}_{a,b}) = \overline{p}^{-1}(a)\)。将 p、I、a、b 分别替换为 -p、-I、-b、-a，即可由此得到关系 \(\mathrm{pr}_{2}(\mathrm{F}_{a,b}) = \overline{p}^{-1}(b)\)。

在 \(c\) 的假设下，\((y,t)\) 属于 \(\overline{p}^{-1}(a)\times\overline{p}^{-1}(c)\)。集合 \(\{y,z\}\) 包含于 \(\overline{p}^{-1}([a,b])\) 的一个连通子集 C 中，集合 \(\{z,t\}\) 包含于 \(\overline{p}^{-1}([b,c])\) 的一个连通子集 C' 中。于是，C \(\cup\) C' 是 \(\overline{p}^{-1}([a,c])\) 的连通子集（TG, I, p.~81, 命题 2），并包含 \(\{y,t\}\)，从而有关系 \((y,t)\in\mathrm{F}_{a,c}\)。

回到定理 3 的证明。映射 p 的每个纤维都是紧的（TG, I, p.~77, 命题 7）。由蒂霍诺夫定理（TG, I, p.~63, 定理 3），积空间 \(K = \prod_{a \in I} \overline{p}^{-1}(a)\) 是紧的。令 \(K'\) 为 K 中的元素 \((y_{a})_{a \in I}\) 组成的集合：它们满足 \(y_{p(x)}\) 等于 x，并且都有 \((y_{a}, y_{b}) \in F_{a,b}\)，对 I 中的每一对 \((a, b)\) 当 a \textless{} b 时成立。定理 3 由下述引理推出。

引理 11. --- a) 集合 K' 非空。

\begin{enumerate}
\def\labelenumi{\alph{enumi})}
\setcounter{enumi}{1}
\tightlist
\item
  设 \((s_{a})_{a\in\mathrm{I}}\) 是 \(K'\) 的一个元素。映射 \(s\colon a\mapsto s_{a}\) 从 I 到 X，是 p 的一个连续截面，且满足 \(s(p(x))=x\)。
\end{enumerate}

对于包含点 \(p(x)\) 的 I 的每个有限子集 S，记 \(K_{S}\) 为 K 中的元素 \((y_{a})_{a\in\mathrm{I}}\) 组成的集合，它们满足关系 \(y_{p(x)} = x\)，并且都有 \((y_{a}, y_{b}) \in \mathrm{F}_{a,b}\)，对 S 中的每一对 \((a, b)\) 当 a \textless{} b 时成立。集合 \(K_{S}\) 在 K 中是闭的（引理 10, a)），并构成 K 的一个递减过滤族，其交为 \(K'\)。

要证明 \(K'\) 非空，只需证明：对于包含点 \(p(x)\) 的 I 的每个有限子集 S，集合 \(K_{S}\) 非空（TG, I, p.~59）。

设 S 是这样的一个子集；将其元素排列成严格递增序列 \((a_{1},\ldots,a_{n})\)，并令 i 为满足 \(p(x)=a_{i}\) 的整数。置 \(y_{a_{i}}=x\)。由引理 10, b)，可以通过归纳构造元素 \(y_{a_{j}}\in\overline{p}^{-1}(a_{j})\)，其中 \(i<j\leqslant n\)，并通过逆向归纳构造元素 \(y_{a_{j}}\in\overline{p}^{-1}(a_{j})\)，其中 \(1\leqslant j<i\)，使得有 \((y_{a_{j}},y_{a_{j+1}})\in\mathrm{F}_{a_{j},a_{j+1}}\)，对每个满足 \(1\leqslant j<n\) 的整数 j 成立。由引理 10, c)，有 \((y_{a},y_{b})\in\mathrm{F}_{a,b}\)，对 S 中的每一对 \((a,b)\) 满足 a\textless b。由于映射 p 是满射，可以对每个 \(a\in I-S\) 选取一个元素 \(y_{a}\in\overline{p}^{-1}(a)\)。如此构造的族 \((y_{a})_{a\in I}\) 属于 \(K_{S}\)，因此 \(K_{S}\) 非空。

我们来证明 \(b\)。根据 \(K'\) 的定义，\(s\) 是 \(p\) 的一个截面，且满足 \(s(p(x)) = x\)。令 \(a \in I\)。我们来证明 \(s\) 在点 \(a\) 处连续。令 U 为 X 中 \(s_a\) 的一个开邻域。由于 \(\overline{p}^{-1}(a)\) 是紧空间（TG, I, p.~77, 命题 7）且全不连通，\(s_a\) 在 \(\overline{p}^{-1}(a)\) 中有一个包含于 U 的开闭邻域 C（TG, II, p.~32, 推论）。集合 C 和 \(\overline{p}^{-1}(a) - C\) 在 \(\overline{p}^{-1}(a)\) 中是闭的，因而是紧的，并且彼此不交。由于 X 是分离的，它们有不交的开邻域 V 和 \(V'\)。集合 \((V \cap U) \cup V'\) 是 X 中纤维 \(\overline{p}^{-1}(a)\) 的一个开邻域；由于映射 \(p\) 是闭的（TG, I, p.~72, 命题 1），\((V \cap U) \cup V'\) 包含形如 \(\overline{p}^{-1}(J)\) 的集合，其中 \(J \subset I\) 是包含 \(a\) 的开区间（I, p.~75, 引理）。置 \(W = V \cap U \cap \overline{p}^{-1}(J)\)。集合 W 在 \(\overline{p}^{-1}(J)\) 中是开的；在 \(\overline{p}^{-1}(J)\) 中也为闭的，因为有 \(\overline{p}^{-1}(J) - W = V' \cap \overline{p}^{-1}(J)\)。令 \(b \in J\)。以 \(a\) 和 \(b\) 为端点的 I 的闭区间包含于 J 中。根据假设，\((s_a, s_b)\) 属于 \(F_{a,b}\)，当 \(a \leqslant b\) 时；\((s_b, s_a)\) 属于 \(F_{b,a}\)，当 \(b \leqslant a\) 时。因此，\(\overline{p}^{-1}(J)\) 中存在包含 \(\{s_a, s_b\}\) 的连通子集。于是，\(s_b\) 属于 \(\overline{p}^{-1}(J)\) 中每个包含 \(s_a\) 的开闭子集，特别地属于 W；从而 \(s_b\) 属于 U。因此有 \(s(J) \subset U\)，这证明了 \(s\) 在点 \(a\) 处连续。

推论. --- 设 X 和 B 是拓扑空间，\(p: X \to B\) 是连续、开、逆紧且分离的映射，其纤维全不连通。设 I 是 R 的一个区间，\(f: I \to B\) 是连续映射，a 是 I 中的一点，x 是 X 中满足 \(f(a) = p(x)\) 的一点。存在连续映射 \(g: I \to X\)，使得 \(p \circ g = f\) 且 \(g(a) = x\)。

置 \(X' = I \times_{B} X\)，并记 \(p'\colon X'\to I\) 和 \(f'\colon X'\to X\) 为典范投影。映射 \(p'\) 是连续、开、逆紧且分离的（I, p.~17, 命题 8 及 p.~27, 命题 4）。由于空间 I 是分离的，空间 \(X'\) 也是分离的（I, p.~26, 注 3）。\(p'\) 的纤维全不连通（I, p.~10, 推论, a)）。由定理 3，存在连续截面 \(s'\)，它是 \(p'\) 的截面，并取值 \((a,x)\) 于 \(a\) 处。从 I 到 X 的映射 \(g = f'\circ s'\) 是连续的；有 \(p\circ g = p\circ f'\circ s' = f\circ p'\circ s' = f\) 且 \(g(a) = x\)，故得推论。

定理 4. --- 设 X 是一个拓扑空间，G 是一个离散群，在 X 上作逆紧作用，p: X → X/G 是典范映射。设 I 是 R 的一个区间，f: I → X/G 是连续映射，a 是 I 中的一点，x 是 X 中满足 \(f(a) = p(x)\) 的一点。存在连续映射 \(\varphi\)：I → X，使得 \(p \circ \varphi = f\) 且 \(\varphi(a) = x\)。

我们先处理 I 是 R 的一个有界闭区间的情形。

根据 TG, III, p.~29, 命题 3，空间 X 是分离的。

令 \(y\) 为 X/G 中的一点，并令 \(x \in X\) 满足 \(y = p(x)\)。\(K_x\) 是 \(x\) 的稳定子，是 G 的有限子群；此外，存在 X 中的邻域 \(U_x\)，其为点 \(x\) 的邻域，以及 X/G 中的邻域 \(V_y\)，其为点 \(y\) 的邻域，使得 \(U_x\) 在 \(K_x\) 下稳定，\(gU_x \cap U_x = \varnothing\)，若 \(g \notin K_x\)，且 \(p\) 诱导从 \(U_x/K_x\) 到 \(V_y\) 的同胚。另外，对每个 \(g \in G\)，\(p\) 诱导从 \(gU_x\) 到 \(V_y\) 的同胚。由于 I 是紧的，存在整数 \(m\) 和 \(n\)，满足 \(m \leqslant 0 \leqslant n\)，以及 I 中元素的有限序列 \((a_i)_{m \leqslant i \leqslant n}\)，使得 \(a_0 = a\)、\(I = {[}a_m,a_n{]}\)，并且对每个 \(i \in \{m, \ldots, n-1\}\)，\(f({[}a_i,a_{i+1}{]})\) 包含于按上述方式构造的 X/G 的一个开集 \(V_{y_i}\) 中。

令 \(x_0\) 为 \(\overline{p}^{-1}(y_0)\) 中满足 \(x \in \mathrm{U}_{x_0}\) 的唯一元素。令 \(q_0\) 为通过取商从 \(\mathrm{U}_{x_0}\) 到 \(\mathrm{U}_{x_0}/\mathrm{K}_{x_0}\) 的典范映射。映射 \(p\) 诱导同胚 \(i_0\)，从 \(\mathrm{U}_{x_0}/\mathrm{K}_{x_0}\) 到 \(\mathrm{V}_{y_0}\)，使得 \(i_0 \circ q_0 = p \mid \mathrm{U}_{x_0}\)。映射 \(q_0\) 是逆紧的（TG, III, p.~29, 命题 3）、开的（TG, I, p.~31, 例 1）且分离的，因为 X 是分离的。它的纤维全不连通，因为它们是有限集。由 III, p.~286 的推论，存在连续映射 \(\varphi_0\colon [a_0,a_1]\to \mathrm{U}_{x_0}\)，使得 \(p\circ \varphi_0 = f\mid [a_0,a_1]\)。

类似地，通过对整数 \(i \in \{0, \ldots, n-1\}\) 归纳，我们构造 \(x_i \in \overline{p}^{-1}(y_i)\) 以及连续映射 \(\varphi_i: [a_i, a_{i+1}] \to X\)，其像包含于 \(\mathrm{U}_{x_i}\) 中，并满足 \(p \circ \varphi_i = f \mid [a_i, a_{i+1}]\)，且满足 \(\varphi_i(a_{i+1}) = \varphi_{i+1}(a_{i+1})\)，当 \(0 \leqslant i < n-1\) 时。

类似地，通过对整数 \(i \in \{m, \ldots, -1\}\) 作逆向归纳，构造 \(x_i \in \overline{p}^{-1}(y_i)\) 以及连续映射 \(\varphi_i: [a_i, a_{i+1}] \to X\)，其像包含于 \(U_{x_i}\) 中，并满足 \(p \circ \varphi_i = f \mid [a_i, a_{i+1}]\)，且满足 \(\varphi_i(a_{i+1}) = \varphi_{i+1}(a_{i+1})\)，当 \(m \leqslant i < 0\) 时。

存在唯一映射 \(\varphi\colon\mathrm{I}\to\mathrm{X}\)，它与 \(\varphi_{i}\) 在 \([a_{i},a_{i+1}]\) 上一致，对满足 \(m\leqslant i<n\) 的 i 成立。它是连续的（TG, I, p.~19, 命题 4），是 \(f\) 到 X 的连续提升，并满足 \(\varphi(a)=x\)。

当 I 紧时，定理得证。在一般情形下，存在序列 \((a_{n})_{n\in\mathbf{N}}\) 和 \((b_{n})_{n\in\mathbf{N}}\)，使得 \((a_{n})\) 以 \(\inf(I)\) 为极限且最终为常值，\((b_{n})\) 以 \(\sup(I)\) 为极限且最终为常值，\(a=a_{0}=b_{0}\)，并且当 I 有最小元（分别有最大元）时，\((a_{n})\)（分别 \((b_{n})\)）为常值。根据上文，对每个整数 \(n\in N\)，存在连续提升 \(\varphi_{n}\) 为 \(f\mid[a_{n},a_{n+1}]\) 到 X 的提升，以及连续提升 \(\varphi_{n}^{\prime}\) 为 \(f\mid[b_{n+1},b_n]\) 到 X 的提升，使得 \(\varphi_{0}(a_{0})=\varphi_{0}^{\prime}(b_{0})=x\)，且 \(\varphi_{n+1}(a_{n+1})=\varphi_{n}(a_{n+1})\)、\(\varphi_{n+1}^{\prime}(b_{n+1})=\varphi_{n}^{\prime}(b_{n+1})\)。存在唯一映射 \(\varphi\colon I\to X\)，它与\(\varphi_{n}\) 在 \([a_{n},a_{n+1}]\) 上一致，并与\(\varphi_{n}^{\prime}\) 在 \([b_{n+1},b_{n}]\) 上一致，对每个 \(n\in N\) 均如此。它是连续的（上引文），是 f 到 X 的连续提升，并满足 \(\varphi(a)=x\)。定理得证。

例. --- 1) 设 X 是一个分离拓扑空间，G 是赋予离散拓扑的有限群，且在 X 上连续作用。于是该作用是逆紧的（TG, III, p.~28, 命题 2）。定理 4 的断言于是直接由定理 3 的推论得到。

\begin{enumerate}
\def\labelenumi{\arabic{enumi})}
\setcounter{enumi}{1}
\tightlist
\item
 设 n 是一个满足 \(\geqslant 0\) 的整数。记 \(P_{n}\) 为次数为 n 的首一多项式 \(P \in C[X]\) 组成的集合，并在其上赋予这样的拓扑：映射 \((c_{0}, \ldots, c_{n-1}) \mapsto X^{n} + c_{n-1}X^{n-1} + \cdots + c_{0}\) 是从 \(C^{n}\) 到 \(P_{n}\) 的同胚。由从 \(C^{n}\) 到 \(P_{n}\) 的映射 p，其中 \(p(z_{1}, \ldots, z_{n}) = (X - z_{1}) \ldots (X - z_{n})\)，是连续的。对称群 \(S_{n}\) 通过置换因子作用在 \(C^{n}\) 上，而 p 通过取商定义了从 \(C^{n}/S_{n}\) 到 \(P_{n}\) 的同胚（TG, VIII, p.~22, 命题 1，I, p.~23, 推论 1 及 TG, VIII, p.~20）。因此得到下述论断：
\end{enumerate}

设 I 是 R 的一个区间，\((c_{0},\ldots,c_{n-1})\) 是从 I 到 C 的连续映射序列，a 是 I 中的一点，\((z_{1},\ldots,z_{n})\) 是一个复数序列，满足 \((\mathrm{X}-z_{1})\ldots(\mathrm{X}-z_{n})=\mathrm{X}^{n}+c_{n-1}(a)\mathrm{X}^{n-1}+\cdots+c_{0}(a)\)。存在一个从 I 到 C 的连续映射序列 \((\lambda_{1},\ldots,\lambda_{n})\)，使得 \(\lambda_{i}(a)=z_{i}\) 对 \(1\leqslant i\leqslant n\) 成立，并且有 \((\mathrm{X}-\lambda_{1}(t))\ldots(\mathrm{X}-\lambda_{n}(t))=\mathrm{X}^{n}+c_{n-1}(t)\mathrm{X}^{n-1}+\cdots+c_{0}(t)\) 对所有 \(t\in I\) 成立。

